\section{Padrão de arquitetura e padrão de projeto}

    \par Durante o presente trabalho de conclusão de curso, faz-se uso de fontes diferentes para o embasamento teórico do tema proposto, porém por se tratar de um tema muito abordado pela literatura ao longo do tempo, nota-se o uso de diferentes termos para referenciar os mesmos objetos de estudo. No caso deste trabalho, destaca-se pelo uso dos termos Padrão de Arquitetura e Padrão de Projeto para denominar o MVC. Nesse sentido cabe uma contextualização para melhor entendimento do objetivo do autor ao usar cada termo, bem como em qual contexto o MVC no presente trabalho se encaixa e qual a sua relação com a Arquitetura Limpa.
    
    \par Pode-se definir um padrão de arquitetura como a organização estrutural fundamental para um sistema de software, definindo seus subsistemas de mais alto nível e as regras que governam suas interações. Um padrão arquitetural estabelece uma estrutura para todos os componentes do sistema, focando na organização global e não em detalhes de implementação \cite{livro:pressman:engenharia:2021}. 
    
    \par O autor da Arquitetura Limpa, Robert C. Martin reforça essa visão ao definir a arquitetura como as decisões que são cruciais para a manutenibilidade e o custo de vida do projeto, como a separação de componentes em camadas independentes \cite{livro:martin:cleanarch}. Portanto, um padrão de arquitetura é o esqueleto do sistema, uma decisão estratégica de alto impacto.

    \par Em contrapartida, um Padrão de Projeto (Design Pattern) oferece uma solução reutilizável para um problema recorrente em um contexto de design mais localizado e de menor escopo dentro do software. Ele descreve o problema, a solução e o momento de aplicar tal solução \cite{livro:pressman:engenharia:2021}. Nessa direção, a obra de \cite{livro:gamma:padroes}, descreve um padrão de projeto como a captura do cerne de uma solução, de forma que possa ser adaptada e usada inúmeras vezes, e tal qual um modelo, pode ser aplicado de maneira diferente de acordo com o contexto \cite{livro:gamma:padroes}.

    \par Robert C. Martin, define o MVC como um padrão de arquitetura, todavia traz a contextualização do mesmo dentro da Arquitetura Limpa, explicando como o MVC é diluído nas camadas da arquitetura limpa, colocando as Views e os Controllers na camada de Adaptadores de Interface e os Models como estruturas de dados, que passam as informações entre as camadas \cite{livro:martin:cleanarch}.

    \par Nesse cenário, Ralph Grove e Eray Ozkan, em seu artigo ``THE MVC-WEB DESIGN PATTERN'', analisam o MVC sob a ótica de um Padrão de Projeto. No entanto, o foco principal do artigo é demonstrar como o padrão evoluiu a partir de sua implementação original no ambiente Smalltalk\footnote{Smalltalk é uma linguagem de programação, concebida por Alan C. Kay no início dos anos 1970 no Palo Alto Research Center (PARC) da Xerox, ela representa uma mudança de paradigma em direção à programação orientada a objetos \cite{livro:bergin:1996:history}.} e se tornou a base arquitetural dos \textit{frameworks} web modernos. Os autores traçam um paralelo detalhado entre as duas versões, mostrando que a natureza do ciclo de requisição-resposta da web forçou uma adaptação significativa no fluxo de controle do padrão MVC.

    \par Para corroborar com a observação teórica de Ralph Grove e Eray Ozkan, Bipin Joshi no livro ``Beginning SOLID Principles and Design Patterns for ASP.NET Developers'', demonstra como o \textit{framework} ASP.NET MVC da Microsoft não apenas utilizava, mas obrigava o uso dessa arquitetura \cite{livro:joshi:2016}.

    \par A literatura apresentada evidencia a dupla natureza do MVC, tratado como Padrão de Projeto por sua reutilização e como Padrão de Arquitetura por seu impacto estrutural. O presente trabalho de conclusão de curso utiliza o \textit{framework} Laravel para demonstrar como os princípios da Arquitetura Limpa podem ser aplicados sobre essa base arquitetural MVC. 
    
    \par Nesse contexto, o MVC é entendido incialmente como um Padrão Arquitetural, mas sob a ótica da Arquitetura Limpa, ele é dividido em componentes. Nessa visão, ele passa a operar como o Padrão de Projeto responsável por toda a camada de interação com o usuário, que é externa ao núcleo do sistema, enquanto a Arquitetura Limpa vai definir a estrutura do software como um todo.